% id: 91-28
% book: 베이직쎈 확률과 통계 · 05 확률변수와 확률분포
% section: 유형 055 연속확률변수의 확률 구하기
% figure: tikz:fig-91-28
% answer: 
% answer_source: 
% variant_level: 0 (원본 전사)
% 이 파일은 build.mjs 가 items.json 에서 만든다. 고칠 때는 items.json 을 고친다.
\smash{\raisebox{4.5mm}{\dmkichul{베이직쎈 확률과 통계 91쪽 28번}}}
\noindent\begin{minipage}[t]{0.58\linewidth}\vspace{0pt}\raggedright $0\le x\le 2$에서 정의된 연속확률변수 $X$의 확률밀도함수 $f(x)$의 그래프가 오른쪽 그림과 같다. $\mathrm{P}(0\le X\le a)=\dfrac{3}{7}$일 때, 상수 $a$의 값을 구하시오. \dmcondfill{$a>\dfrac{1}{2}$}{}\end{minipage}\hfill\begin{minipage}[t]{0.4\linewidth}\vspace{0pt}\centering \resizebox{\ifdim\width>\linewidth\linewidth\else\width\fi}{!}{% 91-28(인쇄 91쪽) — 0≤x≤2 에서 정의된 확률밀도함수 y=f(x) 의 그래프: (0,0)에서 (1/2,4/7)까지 오르고 (2,4/7)까지 수평(안내 파선).
% 스캔 화질이 낮아 TikZ 로 다시 그림. 라벨은 원문에 있는 것만(4/7, O, 1/2, 2, x, y, y=f(x)) · 답(a의 값)은 그림에 넣지 않는다.
\begin{tikzpicture}[x=1.05cm, y=1.3cm, line width=0.6pt]
  \draw[->, >=stealth, line width=0.7pt] (-0.35,0) -- (2.55,0) node[below=1pt, font=\small] {$x$};
  \draw[->, >=stealth, line width=0.7pt] (0,-0.38) -- (0,1.32) node[left=1pt, font=\small] {$y$};
  \draw[densely dotted, line width=0.4pt] (0,1) -- (0.5,1);
  \draw[densely dotted, line width=0.4pt] (0.5,0) -- (0.5,1);
  \draw[densely dotted, line width=0.4pt] (2,0) -- (2,1);
  \draw (0,0) -- (0.5,1) -- (2,1);
  \node[left=2pt, font=\small] at (0,1) {$\dfrac{4}{7}$};
  \node[below=1pt, font=\small] at (0.5,0) {$\dfrac{1}{2}$};
  \node[below=1pt, font=\small] at (2,0) {$2$};
  \node[below left=-1pt and -1pt, font=\small] at (0,0) {$\mathrm{O}$};
  \node[anchor=west, font=\small] at (0.85,1.18) {$y=f(x)$};
\end{tikzpicture}
}\end{minipage}\par
